Everything is computed in Gauge Theory Gravity (Definition \claimref{obj.synthesis.gtg-flat-spacetime-gauge-fields}; Fact \claimref{synthesis.imported-gtg-dictionary}): the curvature is the store-sign operator $\mathcal{R}(B) = -R(\hb(B))$ on physical bivectors, with Ricci map $\mathcal{R}(a) = \partial_{b}\cdot\mathcal{R}(b\wedge a)$, scalar $\mathcal{R} = \partial_{a}\cdot\mathcal{R}(a)$, Einstein map $\mathcal{G}(a) = \mathcal{R}(a) - \tfrac12 a\mathcal{R}$, and Schouten map $\Sch(a) = \tfrac12(\mathcal{R}(a) - \tfrac{\mathcal{R}}{6}a)$; $I$ is the pseudoscalar of the physical frame, $I^{2} = -1$.

\emph{(i)} The Lorentz block of the field strength of the Fact is, as an operator, $F(B) = \mathcal{R}(B) - \tfrac{\Lambda}{3}B$ up to its overall sign ($-\mathcal{F}$ of the dictionary, row 13; it vanishes for the de Sitter field). The Fact's second Euler-Lagrange equation is the frame variation, a trivector-valued linear function of the frame direction $a$:
\[
E(a) \;=\; 2\,\partial_{b}\wedge F\bigl(I(a\wedge b)\bigr),
\]
whose component form is the three-form $\epsilon_{abcd}\,e^{b}\wedge F^{cd}$ with $e^{b}\wedge e^{e}\wedge e^{f}$ read as the trivector $e^{b}\wedge e^{e}\wedge e^{f}$ and $\epsilon_{abcd}$ from $e_{a}\wedge e_{b}\wedge e_{c}\wedge e_{d} = \epsilon_{abcd}I$ (the check computes both on three rational curvature operators at $\Lambda \in \{0, \tfrac37, -2\}$ and finds them equal). Two GA identities evaluate it. For the identity operator, $\partial_{b}\wedge(I(a\wedge b)) = -3\,aI$ (in four dimensions the dual of the bivectors through $a$, wedged back with the frame, returns three copies of the dual of $a$). For every curvature operator,
\[
\partial_{b}\wedge\mathcal{R}\bigl(I(a\wedge b)\bigr) \;=\; \mathcal{G}(a)\,I .
\]
Both sides are linear in $\mathcal{R}$, so it suffices to check it on the two parts of $\mathcal{R} = \mathcal{W} + \Sch\wedge\mathrm{id}$, where $(\Sch\wedge\mathrm{id})(a\wedge b) = \Sch(a)\wedge b + a\wedge\Sch(b)$. On the Weyl part the duality $\mathcal{W}(IB) = I\mathcal{W}(B)$ (dictionary row 5) gives $\partial_{b}\wedge\mathcal{W}(I(a\wedge b)) = \partial_{b}\wedge(I\,\mathcal{W}(a\wedge b))$, which is $I$ times the contraction $\partial_{b}\cdot\mathcal{W}(a\wedge b)$ up to sign, zero because $\mathcal{W}$ is trace-free, while its Einstein map is zero. On the Kulkarni-Nomizu part both sides are linear in the symmetric map $\Sch$, and the check verifies the identity on the ten basis symmetric maps, which span. Hence
\[
E(a) = 2\,\mathcal{G}(a)I - \tfrac{2\Lambda}{3}(-3\,aI) = 2\bigl(\mathcal{G}(a) + \Lambda a\bigr)I .
\]
For the de Sitter field ($\mathcal{R}(B) = \tfrac{\Lambda}{3}B$) $F = 0$ and $E = 0$, and $E \neq 0$ at a wrong $\Lambda$; the check confirms this on the de Sitter gauge field in the flat slicing, the Weyl duality and $\partial_{b}\wedge\mathcal{W}(I(a\wedge b)) = 0$ on the Weyl operators of three gauge fields (Schwarzschild, a Kasner vacuum, a non-vacuum synchronous field), and $\partial_{b}\wedge(I(a\wedge b)) = -3aI$.

\emph{(ii)} By (i), $\mathcal{R} \mapsto E$ factors through the Einstein map, which vanishes on the Weyl part and is injective on the Ricci part (in four dimensions $\mathcal{G}$ determines $\mathcal{R}(a) = \mathcal{G}(a) - \tfrac12 a\,\partial_{b}\cdot\mathcal{G}(b)$, and the Ricci part $\Sch\wedge\mathrm{id}$ is determined by the Ricci map). The algebraic curvature operators form a twenty-dimensional space, the Weyl operators ten of it and the Ricci part ten; so the map has rank $10$ and kernel exactly the Weyl operators, and the nine trace-free components of the Ricci map enter as nine independent conditions. The GTG check finds rank $10$ on the images of the ten Kulkarni-Nomizu operators; the earlier check computes the rank on a spanning set of the twenty-dimensional space in components. A vacuum stationary point has $\mathcal{G}(a) + \Lambda a = 0$, that is $\mathcal{R}(a) = \Lambda a$.

\emph{(iii)} A pure dilation $\delta\hb^{-1}(a) = \delta w\,\hb^{-1}(a)$ (the conformal class fixed) changes the functional by $\int\delta w\,\partial_{a}\wedge E(a)$, so the restricted Euler-Lagrange equation is the pseudoscalar $\partial_{a}\wedge E(a)$. In four dimensions $a\wedge(vI) = (a\cdot v)I$ for vectors, so $\partial_{a}\wedge E(a) = 2\,\partial_{a}\cdot(\mathcal{G}(a) + \Lambda a)\,I = 2(4\Lambda - \mathcal{R})I$, the trace of $\mathcal{G}$ being $-\mathcal{R}$. The GTG check confirms this on the FLRW and Schwarzschild gauge fields at $\Lambda = 0, \tfrac53$, and the earlier check that the map has rank $1$ on the curvature. This is the trace $\mathcal{R} = 4\Lambda$, the single Euler-Lagrange equation of the Einstein-Hilbert action restricted to a conformal class in Theorem \claimref{synthesis.positive-structure-and-natural-scale-actions-do-not-select-the-almost-einstein-scale}(ii) (the Einstein-Cartan term of $S_{\mathrm{MM}}$ is that action and the Euler term is topological).

\emph{(iv)} The coupling to matter is the standard one for the Einstein-Cartan functional: the frame variation of $S_{m}$ is the energy-momentum three-form, and with the normalisation stated the frame equation becomes $\mathcal{G}(a) + \Lambda a + 8\pi G\,T(a) = 0$, $T$ a symmetric linear map on physical vectors (Fact \claimref{synthesis.imported-einstein-cartan-matter-coupling}(a): Kibble 1961, Sciama 1964, Hehl, von der Heyde, Kerlick and Nester 1976; in Gauge Theory Gravity Lasenby, Doran and Gull 1998; the sign is the store's, matching clause (7)). The rest is algebra, checked exactly with $k = 8\pi G$ as a rational symbol on fixed rational $T$ and several $(\Lambda, k)$: given $T$, define $\Ric$ by clause (7); then $\mathcal{G} + \Lambda g + kT = 0$ identically, and conversely, since the map $\Ric \mapsto \mathcal{G}$ is invertible in four dimensions. From clause (7), $\mathcal{R} = 4\Lambda + 8\pi G\,\tau$ and, with the Schouten tensor of Fact \claimref{synthesis.imported-gtg-dictionary}, $\Sch = \tfrac12(\Ric - \tfrac{\mathcal{R}}{6}g) = \tfrac{\Lambda}{6}g - 4\pi G(T - \tfrac{\tau}{3}g)$; its trace-free part is $-4\pi G\,T^{\circ}$ and $J = \tfrac{\mathcal{R}}{6}$ gives $6\cdot\tfrac{J}{4} - 2\pi G\,\tau = \Lambda$. These are the middle slot $-\Sch^{\circ}$ and the relation $\Lambda = 6\rho_{I'} - 2\pi G\,\tau$ of Theorem \claimref{synthesis.infinity-vacuum-incidence-is-a-property-of-the-scale}(4), whose $\nabla_{a}I' = (0, -\Sch^{\circ}_{ab}, \tfrac14\nabla_{a}J)$ then equals $4\pi G(0, T^{\circ}_{ab}, \tfrac{1}{12}\nabla_{a}\tau)$. In vacuum $\Sch = \tfrac{\Lambda}{6}g$ exactly (checked), so $\Sch^{\circ} = 0$ and $J$ is constant, and by Theorem \claimref{synthesis.infinity-vacuum-incidence-is-a-property-of-the-scale}(3) the scale tractor is parallel, which is the own-ray incidence of the infinity vacua of Theorem \claimref{synthesis.einstein-condition-is-infinity-vacuum-incidence}. The value $\tfrac{\Lambda}{6}g$ is the Schouten slot of the Cartan connection $\Omega(a) - \hb^{-1}(a)\wedge V/\ell^{2}$, whose Lorentz curvature block vanishes exactly on the de Sitter field (Theorem \claimref{dark-energy.frozen-point-makes-stelle-west-vector}(ii)).

\emph{Reading.} The field space of $S_{\mathrm{MM}}$ is the pair $(\hb, \Omega)$ (Fact \claimref{dark-energy.imported-macdowell-mansouri-theorem-stelle-west}); in the frozen-point form the constraints $Q(I) = \Lambda/3$, $\ip{V}{I} = 0$, $Q(V) = -3/\Lambda$ fix $\Lambda$ and do not restrict $(\hb, \Omega)$ (Theorem \claimref{dark-energy.frozen-point-makes-stelle-west-vector}(iii)). The functional is not a Lagrange-multiplier or nonnegative-functional rewording of the parallel-tractor condition: its Euler-Lagrange equations are (i) to (iv), and the matter law has $T$ from $S_{m}$. The comparison with the Clausius route uses Theorem \claimref{synthesis.series-wald-entropy-changes-without-flux-on-local-horizons}: on a local edge the Jacobson-Myers density of the series' action is $-\tfrac{3}{4\Lambda}$ times the Cartan curvature of the edge plane, curvature-dependent, while (Loc) asks for a curvature-independent density with no Euler term; so (Loc) is not the Wald entropy of $S_{\mathrm{MM}}$, though both routes give the same field equations. Literature from memory: MacDowell and Mansouri 1977 and Stelle and West 1980 as in the Fact; the Einstein-Cartan coupling is the imported Fact cited in (iv); the Yamabe reading as in the stored no-go.

\emph{In components} (a labelled reading). With $F^{cd} = \tfrac12 F^{cd}{}_{ef}\,e^{e}\wedge e^{f}$ and $F_{abcd} = R_{abcd} - \tfrac{\Lambda}{3}(g_{ac}g_{bd} - g_{ad}g_{bc})$ (Theorem \claimref{synthesis.series-wald-entropy-changes-without-flux-on-local-horizons}), the frame equation is $E_{a} = \tfrac12\epsilon_{abcd}F^{cd}{}_{ef}\,e^{b}\wedge e^{e}\wedge e^{f}$; the contraction of $\epsilon$ with the curvature is the one that gives the Einstein tensor, $E_{a} = 2(\mathcal{G}_{a}{}^{h} + \Lambda\delta_{a}^{h})e_{h}I$. The earlier checks evaluate this on rational curvature tensors (sums of Kulkarni-Nomizu products of rational symmetric matrices), on all fifty-five Kulkarni-Nomizu products of the ten basis symmetric matrices (a spanning set of rank $20$), and the ranks $10$, $9$ (trace-free) and $1$ (dilations), and the matter algebra of (iv) with $k = 8\pi G$ as a rational symbol.
